\section{A tree of daisy cubes with exact positioned supports}
\label{sec:maxmin-chordal-downset-tree}

Simplicial-cube removal supplies a sequence of attachments. A tree description
retains more of their simultaneous geometry: its pieces and overlaps are
daisy cubes, and the same tree records every support in every
coordinate face. This supplies the native pieces and interfaces. Choosing
compatible maximum parents for their union is still a separate problem.
The cube labels are equivalently the extentures of the complement of $H$.
Their fixed literals and free coordinates are both retained in the certificate.

A positioned cube $C$ fixes some coordinates and lets the coordinates in
$d(C)$ vary freely. A coordinate face $P$ is also a positioned cube, including
a vertex or the entire ambient cube. For a finite tree $T$ and a cube label
$C$, let $T_C$ denote the nodes at which that label occurs, together with
the tree edges between them. Unions of such subtrees below are unions of
\emph{vertices and edges}: adjacent vertices do not supply an edge unless
one of the subtrees contains that edge.

\begin{theorem}[A positioned tree certificate for cube-chordal ample classes]
\label{thm:maxmin-chordal-downset-tree}
Let $\varnothing\ne H\subseteq\Q^X$. The following are equivalent.
\begin{enumerate}
\item $H$ is ample and its complete maximal-cube intersection graph
      $\Gamma(H)$ is chordal.
\item There is a finite inclusion-antichain $\mathcal C$ of positioned cubes
      whose union is $H$, and a tree $T$ with nonempty bags
      $\mathcal C_t\subseteq\mathcal C$, such that:
      \begin{enumerate}
      \item every label occurs, every bag has a common word $z_t$, and
            each $T_C$ is connected;
      \item any two intersecting labels occur together in a bag;
      \item for every coordinate face $P$ and $S\subseteq d(P)$,
            the following subcomplex of $T$ is connected or empty:
            \[
             T_{P,S}=\bigcup_{\substack{C\in\mathcal C:\
                    C\cap P\ne\varnothing,\ S\subseteq d(C)}}T_C.
            \]
      \end{enumerate}
\end{enumerate}
In item~2, $\mathcal C$ is automatically the complete maximal-cube list
of $H$. Put
\[
 D_t=\bigcup_{C\in\mathcal C_t}C,\qquad
 F_{tu}=\bigcup_{C\in\mathcal C_t\cap\mathcal C_u}C\quad(tu\in E(T)).
\]
The tree may be chosen so its bags are the maximal cliques of $\Gamma(H)$.
Its pieces $D_t$ are daisy cubes rooted at $z_t$.
Each overlap $F_{tu}$ is nonempty and a daisy cube rooted at both
$z_t,z_u$. If $J=\operatorname{Diff}(z_t,z_u)$, it factors as a full
$J$-cube times a daisy cube on the remaining coordinates.

For every nonempty connected subtree $U\subseteq T$, the union
$H_U=\bigcup_{t\in U}D_t$ is ample and has exactly the cube labels occurring
in $U$ as its maximal cubes. For an edge $tu$ splitting $U$ into $U_1,U_2$,
and every coordinate face $P$,
\begin{align}
 H_{U_1}\cap H_{U_2}&=F_{tu},\label{eq:maxmin-tree-exact-overlap}\\
 \Sh(H_{U_1}\cap P)\cap\Sh(H_{U_2}\cap P)
     &=\Sh(F_{tu}\cap P).
 \label{eq:maxmin-tree-positioned-shadows}
\end{align}
Empty traces have empty shattered families. Every prescribed positive
canonical array realizes each nonempty piece and overlap. In particular,
each has an actual maximum parent by the existing geometric conversion.
If $V\subseteq U$ is a nonempty connected subtree, $H_U$ peels to $H_V$.
Every supplied ordinary map on $H_V$ extends to $H_U$, with all boundary
edges directed into $H_V$ and acyclic orientation outside it. The supplied
map itself may be cyclic. The same assertions hold with an edge overlap
$F_{tu}$ as retained target.
\end{theorem}

\begin{proof}
We first explain the graph tree and the positioned-support condition.
Then a counting identity proves the converse, including that no additional
maximal cube has been silently omitted. Finally we recover all subunions
and their overlaps.

\emph{The graph tree.} A connected chordal graph has a perfect elimination
order $v_1,\ldots,v_m$: the later neighbours $N_i$ of $v_i$ form a clique.
This follows by repeatedly deleting a simplicial vertex, using
\href{../dismantlable-ample/sections/classification/chordal_cube_retention.tex}
{Lemma D.1 of the corner-peeling companion}.
Deleting a simplicial vertex preserves connectivity of a graph with more
than one vertex: replace a path through it by the edge between its two
neighbours. Thus $N_i\ne\varnothing$ for $i<m$.

Take bags $B_i=\{v_i\}\cup N_i$, and join $i$ to the least index $p(i)$
in $N_i$. These edges form a tree rooted at $m$. Since $N_i$ is a clique,
$N_i\subseteq B_{p(i)}$. If $v_j\in B_i$ with $i<j$, this inclusion
carries $v_j$ along the increasing parent path until node $j$.
Hence every graph vertex has a connected occurrence subtree. Every graph
edge is covered by a bag, and bags are cliques. Adjacent contained bags
can be contracted into the larger bag. This preserves the tree properties.
If a remaining bag were properly contained in another, every vertex of
the smaller bag would occur along the path to the larger one, so its first
neighbour would also contain it. Thus after all such contractions the
remaining bags are exactly the maximal cliques.

For item~1 use the vertices of $\Gamma(H)$ as cube labels.
The graph is connected: an edge of the connected ample class lies in a
maximal cube, so a vertex path connects any two such cubes through
intersecting maximal cubes. Cube Helly supplies a common word in each
clique bag: their fixed literals are pairwise consistent.
This establishes conditions~2(a,b).

\emph{All positioned carriers use this same tree.} Fix $P,S$ as in~2(c).
The graph induced on labels with $C\cap P\ne\varnothing$ and
$S\subseteq d(C)$ is connected when nonempty.
Indeed $H\cap P$ is ample. Its positioned $S$-cubes form an ample
anchor family after the $S$-coordinates are removed, by the established
\href{../dismantlable-ample/sections/classification/hall_examples.tex}
{projected-cube carrier lemma} in the companion~\cite{Damp}.
That family is connected; for $S=\varnothing$ use $H\cap P$ itself.
Adjacent anchors lift to a full $(S\cup\{i\})$-cube in $H\cap P$.
Choose a maximal cube of $H$ containing each such lift. Consecutive chosen
cubes share the intervening $S$-cube. Connect each endpoint label to a
chosen $S$-cube within its trace. This gives the required connected graph.
Two adjacent labels in that graph have intersecting occurrence subtrees,
because some bag contains both. Their union of subtrees is therefore
connected, proving condition~2(c). Carrier ampleness is reused here;
the additional conclusion is one tree valid simultaneously for every
position and support.

\emph{The counting converse.} Assume item~2. Two cubes occurring together
intersect at the bag root, so label intersections are exactly intersections
of their occurrence subtrees. The label graph is chordal. To see this
directly, repeatedly remove a leaf of $T$. Labels occurring only at that
leaf have all their neighbours in its clique bag and can be eliminated
as simplicial vertices. Remove the leaf from all other occurrence subtrees
and continue. At the last node the remaining labels form a clique.
This produces a perfect elimination order. Such an order excludes an
induced cycle of length at least four: its earliest vertex on that cycle
would have two nonadjacent later neighbours.

All labels containing any fixed word are pairwise intersecting, so their
occurrence subtrees intersect pairwise. Connected subtrees of a tree have
the Helly property. For completeness, root the tree, and take among the
subtrees the one whose nearest node to the root is farthest from it.
Pairwise intersections force that nearest node to lie in every subtree,
by uniqueness of tree paths. Thus the bags containing the word form a
connected subtree. Consequently, for every coordinate face $P$,
\begin{equation}
 |H\cap P|=\sum_t|D_t\cap P|-\sum_{tu\in E(T)}|F_{tu}\cap P|.
 \label{eq:maxmin-tree-word-count}
\end{equation}
Here $D_t\cap D_u=F_{tu}$: if a cube in bag $t$ intersects one in bag
$u$, their occurrence subtrees meet somewhere, and the path to that
meeting node carries at least one of the two labels across edge $tu$.
Every word in their intersection is therefore in a shared label.

Each bag union is a daisy cube rooted at its common word. Its maximal
cubes are exactly its labels: the top word of any cube relative to that
root belongs to a label, and its entire lower cube belongs to the same
label. A shared-label union has the same property at both roots.
Their traces on $P$ are daisy cubes or empty, hence ample.
Moreover $S\subseteq d(P)$ is shattered by $D_t\cap P$ exactly when
some bag label meets $P$ and has $S\subseteq d(C)$.
Condition~2(c) says that these support occurrences count a connected
tree, with an edge present exactly when a shared label supplies that
support. It follows that
\begin{equation}
 \left|\bigcup_t\Sh(D_t\cap P)\right|
   =\sum_t|\Sh(D_t\cap P)|
            -\sum_{tu\in E(T)}|\Sh(F_{tu}\cap P)|.
 \label{eq:maxmin-tree-support-count}
\end{equation}
Equations~\eqref{eq:maxmin-tree-word-count} and
\eqref{eq:maxmin-tree-support-count}, and ampleness of the individual
traces, give
\[
 |H\cap P|=\left|\bigcup_t\Sh(D_t\cap P)\right|.
\]
Every support on the right has a positioned cube in one trace, so is
strongly shattered by $H\cap P$. The lower Sandwich inequality forces
equality and ampleness of $H\cap P$, with precisely that shattered family.
Equality in either Sandwich inequality characterizes ampleness, as recalled
in the shared ample-class setting \cite{Methods}.

If $P$ is a cube contained in $H$, its entire free support is shattered
by $H\cap P=P$. The last equality puts it in some $\Sh(D_t\cap P)$,
which means $P\subseteq D_t$. The bag downset then puts $P$ inside a
label. Thus there are no new maximal cubes of $H$, and the global label
antichain is its exact maximal-cube list. Its already proved chordal
intersection graph proves item~1.

\emph{Subunions and interfaces.} Restricting to a connected subtree $U$
preserves connected label and carrier subtrees by intersection with $U$.
It also preserves edge coverage: two intersecting label subtrees meeting
$U$ have a common node in $U$, by the Helly property applied to those
two subtrees and $U$. The converse just proved therefore applies to every
connected subunion, with its exact maximal-cube list.

For an edge split, intersecting labels from opposite sides either include
a separator label or meet in a bag on one side; the occurrence path then
makes the label from the other side a separator label. This proves
\eqref{eq:maxmin-tree-exact-overlap}. The connected carrier $T_{P,S}$
meets both sides exactly when it includes the cut edge, exactly when a
shared label supplies an $S$-cube meeting $P$.
The shattered-family equality for subunions proves
\eqref{eq:maxmin-tree-positioned-shadows}.

Shared labels contain both bag roots. Every coordinate where the roots
differ is consequently free in every shared label. Factoring these
coordinates out gives the asserted full cube times daisy cube.
Nonempty interfaces follow in the forward construction since adjacent
clique bags in a connected graph have a shared label; in the converse,
condition~2(c) for the whole cube and empty support makes their union
cover every tree edge. Finally Proposition~\ref{prop:maxmin-all-arrays-cube-gluing}
realizes every positive canonical array on each daisy cube, and
the existing geometric conversion supplies its maximum parent.

To retain a connected subtree $V$, remove a leaf bag outside $V$.
Each of its labels absent from its neighbouring bag occurs only at that
leaf, so all its remaining neighbours in the cube graph lie in the leaf
clique. Remove its exclusive vertices by
Lemma~\ref{lem:maxmin-simplicial-cube-remainder}. This preserves every
retained label and gives an ample remainder with exactly the other labels.
After all exclusive labels are removed, discard the leaf bag and repeat.
A leaf with no exclusive label is discarded without deleting any word.
The procedure ends at $H_V$, proving the relative peeling assertion.
To retain $F_{tu}$, subdivide edge $tu$ by the shared-label bag.
Every occurrence and carrier subtree meeting that edge acquires precisely
its subdivision, so the certificate remains valid; retain the new single
node. Finally the existing relative-peeling extension theorem
\href{../dismantlable-ample/sections/classification/uso_dictionary.tex}
{for ordinary maps} extends every prescribed target map by adding the
deleted corners back as sources. This preserves the target, gives inward
boundary, and places every directed cycle inside the target.
\end{proof}

\begin{remark}[Why positioned support connectivity is required]
\label{rem:maxmin-tree-nonample-path}
In coordinate order $(x_0,x_1,x_2)$, consider the induced path
\[
 000-100-110-111-011.
\]
Its maximal cubes are its four edges, so their intersection graph is a
chordal path. The two end edges both use coordinate $x_0$.
Their support carrier is disconnected. The family has five words but
six shattered supports: the empty set, the three singletons and the
pairs $\{0,1\},\{0,2\}$. Thus it is not ample.
The clique bags and their overlaps are individually daisy cubes;
condition~2(c) detects the inconsistent reuse of the coordinate direction.
The displayed embedding of the path is not isometric. The reconstruction
theorem forces isometry in the prescribed coordinates through ampleness.
\end{remark}

\begin{remark}[The $27$-word control]
\label{rem:maxmin-chordal27-downset-tree}
Use the cube indices in the proof of
Proposition~\ref{prop:maxmin-chordal-unit-array-obstruction}.
One maximal-clique tree is the path $1-0-3-5$ with the following bags:
\begin{center}
\begin{tabular}{cclc}
Node & Root & Cube indices & Piece order\\\hline
$1$ & $4$ & $1,2,4,6$ & $17$\\
$0$ & $0$ & $0,1,2,3,4,5$ & $23$\\
$3$ & $32$ & $2,3,4,5,7$ & $21$\\
$5$ & $40$ & $3,4,7,8,9$ & $16$
\end{tabular}
\end{center}
The three overlaps have orders $16,20,14$; none is a full cube.
In particular an overlap of order $16$ need not be a $4$-cube.
The counting formula gives $17+23+21+16-16-20-14=27$.
All pieces and overlaps admit every positive canonical array, while the
whole class fails the unit array. Thus even on this tree the local
all-array realizations do not justify copying one global array.
The saved certificate and independent verifier check all $729$ coordinate
faces, all $4096$ positioned support carriers, every cut identity, and all
ten connected bag subunions of this example.
\end{remark}

This is a native reconstruction theorem, with exact positioned cubes and
all face shadows. The separate existence of maximum parents for the charts
and overlaps does not choose compatible parents of their union. The remaining
operation is a leaf attachment with these exact face identities and with
coefficient choices free. The $26\to27$ common-unit-guard obstruction of
Proposition~\ref{prop:maxmin-chordal-unit-array-obstruction} remains a
required control for such a construction. No complete forbidden basis or
unrestricted chordal maximum-ancestry theorem is asserted here.

\subsection{Extentures direct the clique tree towards an opposite root}
\label{sec:maxmin-extenture-tree-directions}

The root values may vary along the tree, but each forbidden pattern imposes
a signed separation law. The law uses the actual positioned-support
carriers, including their edges. In particular it applies independently of
whether the class belongs to the peripheral family.

\begin{proposition}[An opposite bag root for every extenture]
\label{prop:maxmin-extenture-opposite-bag-root}
Let $H$ be nonempty and ample with chordal $\Gamma(H)$, and choose a tree
certificate from Theorem~\ref{thm:maxmin-chordal-downset-tree}. Choose any
common word $z_t$ in each bag. For an extenture $R=(A,a)$ of $H$, set
\[
 U_i(R)=\bigcup_{C:\ A\setminus\{i\}\subseteq d(C)}T_C
 \quad(i\in A),\qquad
 M_t(R)=\{i\in A:(z_t)_i=a_i\}.
\]
Then each $U_i(R)$ is a nonempty connected subtree. If $M_t(R)$ is nonempty,
all the subtrees $U_i(R)$, $i\in M_t(R)$, lie in the same component of
$T\setminus\{t\}$. Direct $t$ towards the neighbour in that component.
No edge receives both directions. Consequently every maximal directed path ends
at a node $u$ with
\[
                     z_u|_A=\mathbf1_A-a.
\]
Thus the chosen bag roots, for \emph{every} choice of those roots, contain
the full opposite trace of every extenture.
\end{proposition}

\begin{proof}
We first justify the support carriers and their signs. The projection of
an ample class onto an extenture support $A$ is the full $A$-cube with
exactly the forbidden pattern $a$ deleted. Indeed minimality supplies all
one-bit neighbours of $a$. The projection and its complement are ample;
the complement is nonempty, connected, and has $a$ as an isolated vertex,
so it is just $\{a\}$. Hence every $A\setminus\{i\}$ is shattered.
A maximal cube freeing these coordinates must fix $i$ opposite to $a_i$:
freeing $i$ as well, or fixing it to $a_i$, would realize $R$.
The positioned carrier condition makes $U_i(R)$ connected and nonempty.
It also proves that $t\notin U_i(R)$ whenever $i\in M_t(R)$.

\emph{All matching literals point to one branch.}
Suppose $i,j\in M_t(R)$ had carriers in different components after $t$
is removed. The carrier for $A\setminus\{i,j\}$ contains both subtrees,
so its connectedness puts $t$ in that carrier. There is therefore a cube
in bag $t$ freeing $A\setminus\{i,j\}$. Since it avoids $R$, it fixes
at least one of $i,j$ opposite to the corresponding forbidden bit.
But it contains $z_t$, which matches both bits. This contradiction proves
the common-branch assertion and makes the direction well defined.

\emph{Opposite directions cannot occur.}
Suppose $t$ points to its neighbour $u$ and $u$ points to $t$. Choose
$i\in M_t(R)$ and $j\in M_u(R)$. The subtree $U_i(R)$ is on the $u$
side of the edge and $U_j(R)$ on the $t$ side. They cannot be the same
subtree, so $i\ne j$. The connected carrier for $A\setminus\{i,j\}$
contains the edge $tu$. Because carriers retain edges, some cube label
$C$ occurring in \emph{both} bags frees $A\setminus\{i,j\}$.
It must fix $i$ or $j$ opposite to $a$. The first possibility contradicts
$z_t\in C$ and $(z_t)_i=a_i$; the second contradicts $z_u\in C$ and
$(z_u)_j=a_j$. No such edge exists.

Every node with $M_t(R)\ne\varnothing$ has an outgoing edge. An infinite
walk in a finite tree would repeat an edge in both directions, which was
just excluded. Following the directions therefore terminates at a node
with $M_u(R)=\varnothing$, giving the required opposite trace. For a
singleton support all bag roots already have its opposite bit, and the
same conclusion holds without any directed edge.
\end{proof}

For any positive coefficients on $R$, the sign vector of such an opposite
root makes its guard strictly negative term by term. Different extentures
may end at different nodes. The proposition supplies signed endpoints;
it does not supply one common coefficient array, a negative value for every
row at every root, or a maximum parent. The numerical compatibility between
rows remains a separate question.

For example, the four-word path $000,100,110,111$ has two maximal-clique
bags, with forced roots $100$ and $110$. Its extentures are
$(x_1,x_2)=(0,1)$, $(x_1,x_3)=(0,1)$ and $(x_2,x_3)=(0,1)$.
For the first pattern the opposite trace occurs at $100$; for the last it
occurs at $110$. Thus even this example needs different destinations for
different extentures.

\subsection{One coordinate can prescribe an arbitrary root pattern}
\label{sec:maxmin-tree-arbitrary-root-patterns}

A tree certificate does not make its coordinate data monotone along the
tree. In particular, one cannot assume that the roots switch value in a
given coordinate across at most one tree edge. The phenomenon already
occurs for median classes with unit separator realizations and known
maximum ancestry; it is not an obstruction to that ancestry. Here median
means that the given one-inclusion graph is isometric and the
coordinatewise majority of any three words belongs to the class.

\begin{proposition}[Arbitrary root patterns on a unique clique tree]
\label{prop:maxmin-arbitrary-clique-root-patterns}
Let $T=(V,E)$ be any finite tree with $m\ge1$ edges, and let
$\varepsilon:V\to\{0,1\}$ be arbitrary. There is an ample median class
$H\subseteq\{0,1\}^{2m+2}$ with $3(m+1)$ words and VC dimension two,
whose maximal-cube intersection graph is chordal, such that:
\begin{enumerate}
\item its maximal cliques have a unique clique tree, canonically
      identified with $T$;
\item the cubes in each clique indexed by $v$ have exactly one common
      word $z_v$;
\item in one distinguished coordinate $i$, one has
      $(z_v)_i=\varepsilon(v)$ for every $v\in V$.
\end{enumerate}
Thus even the distinguished coordinate can change on every edge of a
path of arbitrary length. Neither choosing different bag roots nor
choosing a different maximal-clique tree removes those changes.
\end{proposition}

\begin{proof}
Choose a root of $T$ and give its $m$ edges distinct coordinate names.
Let $c_v\in\{0,1\}^E$ indicate the edges on the root-to-$v$ path.
These words form an isometric copy $C_T$ of the tree: the coordinates
where two labels differ are precisely the edges of their joining path.
It is a maximum rank-one class on $E$.
Add a coordinate $i$ and one fresh coordinate $p_v$ for each vertex.
With every unmentioned pendant coordinate zero, put
\[
 H=(C_T\times\{0,1\}_i\times\{0\}^{V})
     \ \cup\ \{(c_v,\varepsilon(v),\mathbf1_{\{v\}}):v\in V\}.
\]
Its graph is the prism $T\mathbin\square K_2$ with one pendant edge
at $(c_v,\varepsilon(v))$ for each $v$. This gives $2m+2$ active
coordinates and $3(m+1)$ distinct words. The prism of a tree is median,
and adding pendant edges preserves medianity. More explicitly, if a
pendant coordinate survives the coordinatewise majority of three words,
that same pendant word occurred at least twice and is the majority.
Otherwise all pendant coordinates vanish; the base majority lies in
$C_T$, and either value of $i$ is present. The constructed graph is
isometric, so this verifies the median assertion.

The only shattered pairs are $\{e,i\}$ with $e\in E$. No base pair
is shattered in the maximum tree, and a pendant coordinate at value
one determines every other coordinate. No triple can be shattered,
since each contains an unshattered pair. Hence
\[
 |\operatorname{Sh}(H)|=1+(2m+2)+m=3(m+1)=|H|,
\]
which proves ampleness and VC dimension two directly.

There is one maximal square $Q_e$ for each base edge $e$, obtained by
multiplying it by $\{0,1\}_i$, and one maximal pendant edge $L_v$ for
each vertex. These are all the maximal cubes. There are no larger
cubes by the preceding pair calculation, and every prism edge belongs
to a square because every vertex of the nontrivial tree has an incident
edge. A pendant edge belongs to no square.
Two squares intersect exactly when their base edges share a vertex.
The edge $L_v$ meets exactly the squares of the base edges incident
with $v$, and distinct pendant edges are disjoint. Consequently the
maximal cliques are exactly
\[
 \mathcal C_v=\{L_v\}\cup\{Q_e:e\text{ is incident with }v\}.
\]
Indeed pairwise incident edges of a tree have a common vertex, so any
clique consisting only of squares is contained in one such clique.
A clique containing $L_v$ is contained in $\mathcal C_v$ as well.
Their clique tree is $T$: each $Q_{uv}$ occurs in precisely the two
bags $\mathcal C_u,\mathcal C_v$, while each $L_v$ occurs only once.
This proves the running-intersection property and hence chordality.
It also proves uniqueness. Any clique tree must join the two bags
containing $Q_{uv}$ directly, since no other bag contains that label.
These $m$ forced edges already form the whole tree on $m+1$ bags.

Finally, the common intersection of $\mathcal C_v$ is the single
prism vertex $z_v=(c_v,\varepsilon(v),0)$. The pendant edge fixes
its base and $i$ coordinates, while any incident square forces its
pendant coordinate to zero. Thus the roots are forced and have the
prescribed distinguished bits. Alternating $\varepsilon$ on a path
proves the last assertion for every length.
\end{proof}

\paragraph{Established membership and explicit checks.}
These classes lie in $\mathcal R_{\rm per}$ by existing operations:
construct the tree by pendant peripheral expansions, duplicate the whole
tree in coordinate $i$, and attach the remaining pendant edges by
singleton peripheral expansions. In particular they are dismantlable.
Their convex realizability and maximum ancestry also follow from the
existing binary-pattern construction
\cite[Theorem~\texttt{thm:two-large-patterns}]{Methods}; no new positive
membership family is being claimed here.

For an explicit certificate, their canonical extentures are the
$\binom m2$ missing pairs of $C_T$, the $(m+1)^2$ pairs
\[
 \{p_v=1,\ x_j=1-(z_v)_j\}
 \quad(v\in V,\ j\in E\sqcup\{i\}),
\]
and the $\binom{m+1}{2}$ pairs $\{p_u=1,p_v=1\}$ for $u\ne v$.
These patterns define exactly the displayed class and are all minimal,
since every coordinate is active. Their number is $q=2m^2+2m+1$.
The old unit-row witness construction is particularly simple here.
For an allowed word, reflect it to zero. A missing pair $01$ imposes
an order on its two positive witness magnitudes. This implication digraph
has no directed cycle: such a cycle would force its active coordinates
to change together, contradicting an isometric path from zero.
Give the coordinates weights $1,\ldots,2m+2$ in topological order.
Their negatives strictly satisfy all the reflected binary inequalities;
a missing $11$ row is automatically negative. Reflect back to obtain
integer witnesses of absolute value at most $2m+2$ and unit-row slack
at least one. The existing bounded-realization conversion gives a
sufficient maximum-parent rank $2m+2$ and at most
$q+2(2m+2)=2m^2+6m+5$ auxiliaries. These bounds are not claimed minimal.

The independent replay checks a path and a branching tree with four
edges, with prescribed root patterns, and the one-edge boundary case.
It reconstructs all cubes, maximal cliques and their forced tree edges,
checks majority closure and the full shattered family, and verifies every
unit-row witness. The statement about arbitrary trees and patterns is
proved symbolically. The additional structural conclusion is unrestricted
coordinate variation on the clique tree; it does not supply compatible
parents for a general chordal union or prove corner-addition closure.
